\lecture{12}{2026-10-27}{Post-Quantum Cryptography and Learning with Errors}

\section*{Quantum Attacks and the Hidden Subgroup Problem}

Classical discrete logarithm algorithms (Pollard's \(\rho\), Pohlig--Hellman, Index Calculus) require exponential or subexponential time.
However, quantum algorithms can solve discrete logarithms in polynomial time.

\begin{definition}[Abelian Hidden Subgroup Problem]
    Let \(G\) be a finitely generated abelian group and \(S\) a finite set.
    Suppose we are given oracle access to a function \(f \colon G \to S\) with the property that there exists a subgroup \(H \le G\) such that
    \[
        f(x) = f(y) \iff x - y \in H \iff x + H = y + H.
    \]
    The \emph{Abelian Hidden Subgroup Problem} (HSP) asks to determine a generating set for the hidden subgroup \(H\).
\end{definition}

Shor's algorithm~\cite{Shor1994} solves the abelian hidden subgroup problem in polynomial time on a quantum computer.

\begin{proposition}[Discrete Logarithm as an Abelian HSP]
    Let \(G = \langle P \rangle\) be a cyclic group of prime order \(q\) (such as a prime-order subgroup of \(E(\mathbb{F}_p)\) or \(\mathbb{Z}_p^\times\)).
    Let \(Q = xP\) for an unknown scalar \(x \in \mathbb{Z}_q\).
    The discrete logarithm problem can be cast as an abelian HSP over \(\mathbb{Z}_q \times \mathbb{Z}_q\).
\end{proposition}

\begin{proof}
    Define \(f \colon \mathbb{Z}_q \times \mathbb{Z}_q \to \langle P \rangle\) by
    \[
        f(a, b) = aP + bQ.
    \]
    Since \(Q = xP\), we have \(f(a, b) = (a + bx)P\).
    Then \(f(a_1, b_1) = f(a_2, b_2)\) if and only if
    \begin{align*}
        (a_1 + b_1 x)P = (a_2 + b_2 x)P
        &\iff a_1 + b_1 x \equiv a_2 + b_2 x \pmod q \\
        &\iff a_1 - a_2 \equiv -x(b_1 - b_2) \pmod q \\
        &\iff (a_1 - a_2, b_1 - b_2) \in \langle (-x, 1) \rangle.
    \end{align*}
    Thus, \(f\) hides the subgroup \(H = \langle (-x, 1) \rangle \le \mathbb{Z}_q \times \mathbb{Z}_q\).
    Finding a non-zero element in \(H\) immediately reveals the secret discrete logarithm \(x\).
\end{proof}

Consequently, Shor's algorithm completely breaks both finite field and elliptic curve Diffie--Hellman key exchange.
This motivates constructing cryptographic primitives whose security relies on problems resistant to quantum attacks.

\section*{A First Attempt: Linear Algebra Diffie--Hellman}

As an alternative to group-based Diffie--Hellman, one might consider using matrix multiplication over a finite field \(\mathbb{Z}_q\).

\paragraph*{Protocol}
Fix public parameters: dimensions \(m, n \in \mathbb{N}\), a prime \(q\), and a uniformly random matrix \(A \in \mathbb{Z}_q^{m \times n}\).
\begin{enumerate}
    \item Alice samples a secret vector \(s_A \in \mathbb{Z}_q^n\), computes \(u_A = A s_A \in \mathbb{Z}_q^m\), and sends \(u_A\) to Bob.
    \item Bob samples a secret vector \(s_B \in \mathbb{Z}_q^m\), computes \(u_B^T = s_B^T A \in \mathbb{Z}_q^{1 \times n}\), and sends \(u_B^T\) to Alice.
    \item Both parties compute the shared secret:
    \[
        s_B^T u_A = s_B^T A s_A = u_B^T s_A \in \mathbb{Z}_q.
    \]
\end{enumerate}

\paragraph*{Insecurity}
This scheme is completely broken.
An eavesdropper observes \(A\) and \(u_A = A s_A\).
Finding Alice's secret \(s_A\) amounts to solving a system of \(m\) linear equations in \(n\) variables over \(\mathbb{Z}_q\).
Gaussian elimination solves this in \(O(m n^2)\) field operations, breaking the scheme in polynomial time classically.

\section*{Learning with Errors (LWE)}

To prevent Gaussian elimination from recovering the secret vectors, Regev~\cite{Regev2005} introduced the idea of perturbing the linear systems with small random noise.

\subsection*{Attempted Key Exchange with Noise}

In the noisy setting:
\begin{itemize}
    \item Alice samples a small secret \(s_A \in \mathbb{Z}_q^n\) and a small error vector \(e_A \in \mathbb{Z}_q^m\), sending
    \[
        b_A = A s_A + e_A.
    \]
    \item Bob samples a small secret \(s_B \in \mathbb{Z}_q^m\) and a small error vector \(e_B \in \mathbb{Z}_q^n\), sending
    \[
        b_B^T = s_B^T A + e_B^T.
    \]
\end{itemize}
Then Alice computes
\[
    s_B^T b_A = s_B^T A s_A + s_B^T e_A
\]
and Bob computes
\[
    b_B^T s_A = s_B^T A s_A + e_B^T s_A.
\]
This does not directly yield a key exchange protocol because the two parties' computed values differ by the error term \(s_B^T e_A - e_B^T s_A\).
However, because this error term is small, we can use this mechanism to construct a public-key message encryption scheme instead.

\subsection*{Public-Key Encryption from LWE}

\paragraph*{Key Generation}
Alice samples a uniformly random matrix \(A \in \mathbb{Z}_q^{m \times n}\), a small secret vector \(s_A \in \mathbb{Z}_q^n\), and a small error vector \(e_A \in \mathbb{Z}_q^m\).
She publishes the public key
\[
    \mathsf{pk} = (A, b_A = A s_A + e_A),
\]
and keeps \(\mathsf{sk} = s_A\) as her secret key.

\paragraph*{Encryption}
To encrypt a single-bit message \(m \in \{0, 1\}\), Bob samples a small secret vector \(s_B \in \mathbb{Z}_q^m\), a small error vector \(e_B \in \mathbb{Z}_q^n\), and a small error scalar \(e'_B \in \mathbb{Z}_q\).
He computes
\[
    c_1 = s_B^T A + e_B^T \in \mathbb{Z}_q^{1 \times n}, \qquad c_2 = s_B^T b_A + e'_B + \lfloor q/2 \rfloor m \in \mathbb{Z}_q,
\]
and sends the ciphertext \(c = (c_1, c_2)\) to Alice.

\paragraph*{Decryption}
Using her secret key \(s_A\), Alice computes
\begin{align*}
    c_2 - c_1 s_A &= \left(s_B^T b_A + e'_B + \lfloor q/2 \rfloor m\right) - \left(s_B^T A + e_B^T\right) s_A \\
    &= s_B^T (A s_A + e_A) + e'_B + \lfloor q/2 \rfloor m - s_B^T A s_A - e_B^T s_A \\
    &= \lfloor q/2 \rfloor m + \underbrace{(s_B^T e_A - e_B^T s_A + e'_B)}_{\text{total error } E} \pmod q.
\end{align*}
Alice recovers the message up to the small error term \(E = s_B^T e_A - e_B^T s_A + e'_B\).
If \(m = 0\), the result is \(E \approx 0 \pmod q\); if \(m = 1\), the result is \(\lfloor q/2 \rfloor + E \approx q/2 \pmod q\).
Alice can thus reliably distinguish between \(m = 0\) and \(m = 1\) by checking which value is closer:
\[
    m' =
    \begin{cases}
        0 & \text{if } |c_2 - c_1 s_A|_q < q/4, \\
        1 & \text{otherwise}.
    \end{cases}
\]
Decryption succeeds correctly whenever the total noise satisfies \(|E|_q < q/4\).

Notice that this basic scheme only encrypts a single bit at a time; in the next lecture we will see how to improve this.

\paragraph*{A caveat: What does ``small'' mean?}
Elements of \(\mathbb{Z}_q\) are modular residues.
To speak of an error being ``small'', we identify elements with their centered representatives in the interval \([-\lfloor q/2 \rfloor, \lfloor q/2 \rfloor]\) and define the centered modular absolute value
\[
    |x|_q = \min_{k \in \mathbb{Z}} |x + kq|.
\]
In practice, error coordinates are sampled from small distributions such as:
\begin{itemize}
    \item \textbf{Uniform bounded:} each coordinate is drawn uniformly from \([-\beta, \beta] \cap \mathbb{Z}\) for \(\beta \ll q\);
    \item \textbf{Centered Binomial Distribution (CBD):} with parameter \(k\), sampled as \(\sum_{j=1}^k a_j - \sum_{j=1}^k b_j\) for independent uniform bits \(a_j, b_j \in \{0, 1\}\);
    \item \textbf{Discrete Gaussian:} where each coordinate is sampled with probability proportional to \(\exp(-|x|_q^2 / (2\sigma^2))\).
\end{itemize}